\section{Wheel-contact constraints}
\label{sec:constraints}

The vehicle has two wheels, each treated as an ideal rigid disk rolling
without slipping on a planar ground. Each contact point yields two scalar
constraints --- one along the rolling direction and one perpendicular to it.
Below we derive both, show that the lateral pair degenerates into a single
body-level constraint, and prove that this constraint is non-holonomic.

\subsection{Wheel positions and contact velocities}

The left wheel sits a distance $\trackw/2$ along $\ybody$ from the body
origin; the right wheel sits the same distance along $-\ybody$:
\begin{align}
    r_L &= \begin{pmatrix} p_x - \half\trackw\sin\theta \\
                           p_y + \half\trackw\cos\theta \end{pmatrix},
    &
    r_R &= \begin{pmatrix} p_x + \half\trackw\sin\theta \\
                           p_y - \half\trackw\cos\theta \end{pmatrix}.
    \label{eq:wheel-positions}
\end{align}
Differentiating in time,
\begin{align}
    \dot r_L &= \begin{pmatrix}
        \dot p_x - \half\trackw\,\dot\theta\cos\theta \\
        \dot p_y - \half\trackw\,\dot\theta\sin\theta
    \end{pmatrix},
    &
    \dot r_R &= \begin{pmatrix}
        \dot p_x + \half\trackw\,\dot\theta\cos\theta \\
        \dot p_y + \half\trackw\,\dot\theta\sin\theta
    \end{pmatrix}.
\end{align}

\subsection{Lateral constraint (no side slip)}

Each wheel cannot move along its axle. Both axles are aligned with $\ybody$,
so for each wheel the constraint is $\dot r_w \cdot \ybody = 0$:
\begin{align}
    \dot r_L \cdot \ybody
    &= -\dot p_x \sin\theta
       + \half\trackw\,\dot\theta\cos\theta\sin\theta
       + \dot p_y \cos\theta
       - \half\trackw\,\dot\theta\sin\theta\cos\theta \notag\\
    &= -\dot p_x \sin\theta + \dot p_y \cos\theta.
\end{align}
The $\dot\theta$-dependent cross terms cancel because the wheel offset is
itself along $\ybody$, so the spin-induced contact velocity is along
$\xbody$, not $\ybody$. The right wheel gives the same expression. Both
lateral constraints therefore collapse into a single body-level Pfaffian
constraint:
\begin{equation}
    \boxed{\;-\sin\theta\;\dot p_x + \cos\theta\;\dot p_y = 0.\;}
    \label{eq:noslip}
\end{equation}

\subsection{Longitudinal constraint (pure rolling)}

The forward velocity of each wheel contact equals its rim velocity:
\begin{align}
    \dot r_L \cdot \xbody &= v_L,
    &
    \dot r_R \cdot \xbody &= v_R.
\end{align}
Expanding,
\begin{align}
    v_L &= \dot p_x \cos\theta + \dot p_y \sin\theta - \half\trackw\,\dot\theta
         = v - \half\trackw\,\omega, \\
    v_R &= \dot p_x \cos\theta + \dot p_y \sin\theta + \half\trackw\,\dot\theta
         = v + \half\trackw\,\omega.
\end{align}
Solving for the body twist,
\begin{equation}
    \boxed{\;
        v = \half(v_L + v_R),
        \qquad
        \omega = \frac{v_R - v_L}{\trackw}.
    \;}
    \label{eq:wheel-to-body}
\end{equation}
This is the map used at the top of \texttt{DiffDriveDynamics::continuous}
in \texttt{core/dynamics.hpp}.

\subsection{Pfaffian form}

Stacking the three constraints,
\begin{equation}
    \underbrace{\begin{bmatrix}
        -\sin\theta & \cos\theta & 0 \\
        \phantom{-}\cos\theta & \sin\theta & -\half\trackw \\
        \phantom{-}\cos\theta & \sin\theta & +\half\trackw
    \end{bmatrix}}_{A(q)}
    \begin{bmatrix}\dot p_x \\ \dot p_y \\ \dot\theta\end{bmatrix}
    \;=\;
    \begin{bmatrix} 0 \\ v_L \\ v_R \end{bmatrix}.
    \label{eq:pfaffian}
\end{equation}
Row~1 is homogeneous in $\dot q$ --- no actuator drives it. This is the
non-holonomic part. Rows~2--3 are inhomogeneous, driven by the wheel
commands.

\subsection{The constraint is non-holonomic}

A Pfaffian constraint $a(q) \cdot \dot q = 0$ is integrable to a function
$h(q) = \mathrm{const}$ if and only if the dual distribution
$\Delta = \ker a(q)$ is involutive (Frobenius theorem). For
\eqref{eq:noslip},
\begin{equation}
    \Delta = \mathrm{span}\{X_1, X_2\},
    \qquad
    X_1 = \begin{pmatrix}\cos\theta \\ \sin\theta \\ 0\end{pmatrix},
    \quad
    X_2 = \begin{pmatrix}0 \\ 0 \\ 1\end{pmatrix},
\end{equation}
where $X_1$ generates pure forward motion and $X_2$ generates pure rotation
in place. The Lie bracket is
\begin{equation}
    [X_1, X_2]
    \;=\; \frac{\partial X_1}{\partial q}\,X_2 \;-\; \frac{\partial X_2}{\partial q}\,X_1
    \;=\; \begin{pmatrix} -\sin\theta \\ \phantom{-}\cos\theta \\ 0 \end{pmatrix} \;-\; 0.
\end{equation}
This vector is precisely the body-lateral direction \eqref{eq:noslip}
forbids, so $[X_1, X_2] \notin \Delta$. The distribution is therefore
non-involutive, \eqref{eq:noslip} is non-integrable, and the system is
non-holonomic. \(\square\)

\begin{remark}[Chow's theorem and parking]
The set $\{X_1,\, X_2,\, [X_1, X_2]\}$ spans $\R^{3}$ at every $q$. By
Chow's theorem the system is locally controllable: any configuration
$(p_x, p_y, \theta)$ can be reached by combining forward motion and
rotation, even though the vehicle has only two control inputs and cannot
translate sideways instantaneously. Parallel parking is the everyday
demonstration of this fact.
\end{remark}

\subsection{Consistency with the implementation}
\label{subsec:constraint-consistency}

The state ODE in \texttt{dynamics.hpp} is written directly in body
coordinates,
\begin{equation}
    \dot p_x = v\cos\theta,
    \qquad
    \dot p_y = v\sin\theta,
\end{equation}
which substituted into \eqref{eq:noslip} gives
$-v\sin\theta\cos\theta + v\sin\theta\cos\theta = 0$ identically. The
constraint is built into the parameterisation; the simulator does not need
projection or Lagrange multipliers to enforce it at runtime. The numerical
verification of this identity over a long rollout will live in
\texttt{analysis/dynamics\_tour.ipynb} (a future ``Test 5: constraint
satisfaction'' cell).
